% Source: Romandini, 08_KeyEscrow slides; lecture of 21 Nov 2026, last 15 minutes of the
%         morning (motivation and the naive split; to be continued in the afternoon).
\section{Secret Sharing and Key Escrow}
\label{sec:secret-sharing}

\subsection{The problem with private keys}

In many situations a secret is known to a single person: the symmetric key that encrypts a file or a whole disk (FileVault on macOS, BitLocker on Windows), a private key, a password or PIN. What happens if that person forgets or loses it, leaves the company, or dies, and someone else legitimately needs the data?

\begin{example}[Cantor Fitzgerald, 2001]
The financial firm lost about 700 employees in the attacks of 11 September 2001. Rebuilding the business started with a search for the passwords that only they had known, so that their accounts and files could be opened.
\end{example}

There is also a legal and political side: governments would like to read encrypted communications under a court order or in case of danger to national security. In 2016 the FBI demanded that Apple help unlock the iPhone of a dead terrorist in San Bernardino; Apple refused, and the FBI eventually paid a third party to break into the phone.

\begin{definition}[Escrow, key escrow]
\textbf{Escrow} is a contractual arrangement in which a third party receives and disburses money or documents for the primary parties, the disbursement depending on agreed conditions (like PayPal holding the payment between an untrusting buyer and seller). In \textbf{key escrow} private keys are held by a third party such that, under certain circumstances, an \emph{authorised} party may access or reveal them.
\end{definition}

\paragraph{Naive solutions.} Give a copy of your key to a trusted figure, e.g.\ the company's Chief Security Officer (CSO); or, rather than handing it over, store it on your own disk encrypted with the CSO's public key, so only the CSO can recover it. Two concerns remain: what if you \emph{do not trust} the CSO, who alone can now read everything? And what if \emph{disaster} strikes the CSO, so that the only recovery path is gone? Both are single points of failure, of confidentiality and of availability respectively. The way out is to split the secret among several people.

\begin{remark}[Escrow in everyday systems]
Beyond the slides. Disk-encryption recovery keys are routinely escrowed: BitLocker keys in the company's directory (Active Directory, Entra ID), FileVault keys at Apple or in the company's device management. It is a good idea, but the escrow must stay reachable when it is needed: in the CrowdStrike outage (\cref{sec:security-basics}) some recovery keys were stored on servers that were themselves not booting.
\end{remark}

\subsection{Secret sharing}

Split the secret into $n$ pieces and give each to a different security officer (or encrypt the $i$-th piece with officer $i$'s public key and keep all of them on your disk), so that:
\begin{itemize}
    \item \textbf{Confidentiality}: no proper subset of the officers can reconstruct the secret any more easily than by brute force;
    \item \textbf{Availability}: all $n$ officers together can reconstruct it.
\end{itemize}

\paragraph{Why not just cut the key into pieces?} The first idea proposed in class: take a 128-bit key and give each colleague a consecutive slice of it, say 10 bits each. Availability holds, since all the slices in order give the key back. Confidentiality fails, because each piece \emph{is} information about the secret. Whoever holds the first 10 bits has to brute-force only the remaining 118, i.e.\ $2^{118}$ keys instead of $2^{128}$. A group of colluders who hold half the slices has to search only the other half. The requirement is stronger: no individual and no unauthorised group may learn \emph{anything} about the secret. The XOR construction below achieves exactly that.

\paragraph{Two officers.} Generate a random string $R$ as long as the secret $S$, and give
\[
    S_1 = R, \qquad S_2 = R \xor S.
\]
Each officer alone has no information about $S$: one holds a random string, the other $S$ encrypted with a one-time pad (\cref{sec:symmetric}), which is information-theoretically perfect. Together they recover
\[
    S_1 \xor S_2 = R \xor (R \xor S) = (R \xor R) \xor S = 0 \xor S = S.
\]

\paragraph{$n$ officers.} For a $k$-bit secret $S$, generate $n-1$ random $k$-bit strings $R_1, \dots, R_{n-1}$ and give
\[
    S_i = R_i \quad (i = 1, \dots, n-1), \qquad S_n = R_1 \xor R_2 \xor \dots \xor R_{n-1} \xor S.
\]
All $n$ shares together give $S_1 \xor \dots \xor S_n = S$, because every $R_i$ appears twice and cancels. Any $n-1$ shares are independent uniformly random strings: they say nothing about $S$.

\subsection{Threshold schemes}

The $n$-way scheme has an availability problem: if one officer dies or loses the share, the remaining shares are useless and the secret is lost. The requirement that \emph{all} shares be present must be relaxed.

\begin{definition}[$(t, n)$-threshold scheme]
For $t \le n$, the secret is split into $n$ \emph{shares} such that any $t$ (or more) shares suffice to recover it, while fewer than $t$ shares give no information about it.
\end{definition}

\begin{definition}[Shamir's secret sharing \cite{shamir1979}]
Choose a public prime $p$ larger than both $S$ and $n$, and a random polynomial of degree $t-1$ whose constant term is the secret:
\[
    g(x) = (a_{t-1} x^{t-1} + a_{t-2} x^{t-2} + \dots + a_1 x + S) \bmod p,
\]
with $a_1, \dots, a_{t-1}$ chosen at random. The shares are $S_i = g(i)$ for $i = 1, \dots, n$; the polynomial (its coefficients and $S$) is then discarded.
\end{definition}

Any $t$ shares give $t$ equations in the $t$ unknowns $a_1, \dots, a_{t-1}, S$, so they determine the polynomial and $S = g(0)$. This is the generalisation of a familiar fact: two points determine a line, three points a parabola, $t$ points a polynomial of degree $t-1$.

\paragraph{Why fewer than $t$ shares reveal nothing.} Take $t = 2$, so $g(x) = (ax + S) \bmod p$ is a line and $S$ is where it crosses the vertical axis. With two shares the line is fixed and so is $S$. With a single share $(i, S_i)$, for \emph{every} candidate secret $S'$ there is exactly one line through $(0, S')$ and $(i, S_i)$: all secrets are equally possible, as with the one-time pad. Working modulo a prime matters: over the ordinary integers a share would leak partial information (e.g.\ bounds on $S$), while in the finite field every value is equally likely.

\begin{example}[A $(2,3)$-threshold scheme, from the slides]
\label{ex:shamir}
Let $n = 3$, $t = 2$, $S = 6$, $p = 17$. A random polynomial of degree $t - 1 = 1$ is $g(x) = (4x + 6) \bmod 17$. The three officers receive
\[
    g(1) = 10, \qquad g(2) = 14, \qquad g(3) = (12 + 6) \bmod 17 = 1.
\]
Officers 1 and 2 know that $g(x) = (ax + S) \bmod 17$, so
\[
    a + S \equiv 10, \qquad 2a + S \equiv 14 \pmod{17}.
\]
Subtracting the first equation from the second gives $a = 4$, hence $S = 10 - 4 = 6$, the original secret. Any other pair works the same way.
\end{example}

\paragraph{Lagrange interpolation.} Beyond the slides: rather than solving the system, the secret can be computed directly from any $t$ shares $(x_1, y_1), \dots, (x_t, y_t)$ as
\[
    S = g(0) = \sum_{i=1}^{t} y_i \prod_{j \neq i} \frac{x_j}{x_j - x_i} \pmod p,
\]
where the division is multiplication by the inverse modulo $p$ (\cref{sec:asymmetric}). A short implementation, with $p = 2^{127} - 1$ (a prime) so that the secret can be a 126-bit key:

\begin{minted}{python}
import secrets

P = 2**127 - 1

def split(secret, t, n):
    coeffs = [secret] + [secrets.randbelow(P) for _ in range(t - 1)]
    def g(x):
        return sum(c * pow(x, i, P) for i, c in enumerate(coeffs)) % P
    return [(x, g(x)) for x in range(1, n + 1)]

def recover(shares):
    secret = 0
    for xi, yi in shares:
        num = den = 1
        for xj, _ in shares:
            if xj != xi:
                num = num * -xj % P
                den = den * (xi - xj) % P
        secret = (secret + yi * num * pow(den, -1, P)) % P
    return secret

key = secrets.randbelow(P)
shares = split(key, t=3, n=5)
print(recover(shares[:3]) == key, recover([shares[0], shares[2], shares[4]]) == key)
print(recover(shares[:2]) == key)
\end{minted}

The first line prints \texttt{True True}: any three of the five shares recover the key. The second prints \texttt{False}: two shares yield a value, but it is just a random number unrelated to the key. Note that the coefficients come from \texttt{secrets}: with a predictable generator, the \enquote{random} polynomial and therefore the secret could be guessed (\cref{sec:base-mechanisms}).

\begin{remark}[Limits of secret sharing]
The shares protect the secret at rest, but to use it someone must reconstruct it, and at that moment the whole secret exists again in one place, which must be trusted. Moreover $t$ colluding officers can always recover it: choosing $t$ is a trade-off between confidentiality (higher $t$) and availability (lower $t$, more spare shares). Beyond the slides, \emph{threshold cryptography} goes further and lets $t$ parties sign or decrypt jointly without ever reconstructing the key.
\end{remark}

\subsection{Escrow and secret sharing in the real world}

\begin{itemize}
    \item \textbf{Clipper chip} (1993): a US government proposal for telephone encryption with each device's key escrowed in two halves by two agencies. It was abandoned after Matt Blaze showed in 1994 that the escrow could be bypassed: the field that let the agencies find the key was protected by only a 16-bit checksum, so a forged field that passed the check could be found by trial and error.
    \item \textbf{DNSSEC root key}: the key that signs the root of the Internet's naming system is used only in public \emph{ceremonies}, in which several trusted community representatives from different countries must be present, each holding a part of what is needed to activate it (e.g.\ smart cards, of which a quorum is required).
    \item \textbf{Shamir's scheme in practice}: splitting backup and recovery keys, e.g.\ the unseal key of HashiCorp Vault (by default 5 shares, any 3 needed) and the backup phrases of cryptocurrency wallets (SLIP-39).
    \item \textbf{Lawful access}: the policy debate on giving governments access to encrypted data is still open. Security experts warn that any escrow mechanism is itself a valuable target and a source of new vulnerabilities, as Clipper showed \cite{abelson2015}.
\end{itemize}

\begin{remark}[Escrow encryption keys, never signing keys]
Beyond the slides. Escrow makes sense for keys that \emph{protect data}: if the key is lost, the data is lost. It makes no sense for \emph{signing} keys: losing one costs nothing (you generate a new pair and get a new certificate), while an escrowed copy means that someone else could sign in your name, which destroys non-repudiation (\cref{sec:signatures}). Well-designed systems therefore use separate key pairs for encryption (escrowed or backed up) and for signatures (never copied).
\end{remark}

\begin{summary}
\begin{itemize}
    \item A secret known to one person is lost with that person. Key escrow lets an authorised third party recover keys under agreed conditions; a single trustee (the CSO) is a single point of failure for both confidentiality and availability.
    \item $n$-way XOR sharing: $S_i = R_i$ for $i < n$, $S_n = R_1 \xor \dots \xor R_{n-1} \xor S$. All shares are needed; any $n-1$ reveal nothing (one-time pad).
    \item $(t, n)$ threshold: any $t$ of $n$ shares recover the secret, fewer reveal nothing. Shamir: random polynomial of degree $t-1$ mod $p$ with constant term $S$, shares $g(1), \dots, g(n)$; $t$ shares determine the polynomial and $S = g(0)$.
    \item Real world: Clipper (escrow bypassed, abandoned), DNSSEC root key ceremonies, Shamir for backup keys (Vault, wallets). Escrow encryption keys, never signing keys.
\end{itemize}
\end{summary}

\subsection*{Exercises}

\begin{exercise}
An 8-bit secret $S = \texttt{10110010}$ is shared among three officers with $R_1 = \texttt{01101100}$ and $R_2 = \texttt{11100001}$. Compute the third share and verify the reconstruction. What do officers 1 and 3 learn about $S$ if they collude?
\end{exercise}
\begin{solution}
$S_3 = R_1 \xor R_2 \xor S = \texttt{01101100} \xor \texttt{11100001} \xor \texttt{10110010} = \texttt{00111111}$. Reconstruction: $S_1 \xor S_2 \xor S_3 = \texttt{10001101} \xor \texttt{00111111} = \texttt{10110010} = S$. Officers 1 and 3 can compute $S_1 \xor S_3 = R_2 \xor S$, which is $S$ encrypted with the one-time pad $R_2$: they learn nothing about $S$.
\end{solution}

\begin{exercise}
In \cref{ex:shamir}, officers 1 and 3 want to recover the secret. Set up and solve the equations.
\end{exercise}
\begin{solution}
$a + S \equiv 10$ and $3a + S \equiv 1 \pmod{17}$. Subtracting: $2a \equiv -9 \equiv 8$, so $a \equiv 4$ (since $2 \cdot 4 = 8$). Then $S = 10 - 4 = 6$.
\end{solution}

\begin{exercise}
A $(3,5)$-threshold scheme modulo $p = 17$ gives officers 1, 2, 3 the shares $g(1) = 1$, $g(2) = 12$, $g(3) = 10$. Recover the secret with Lagrange interpolation. Why could officers 1 and 2 alone not do it?
\end{exercise}
\begin{solution}
With $x_1, x_2, x_3 = 1, 2, 3$ the weights $\prod_{j \neq i} x_j / (x_j - x_i)$ are
$\frac{2 \cdot 3}{(2-1)(3-1)} = 3$, $\frac{1 \cdot 3}{(1-2)(3-2)} = -3$, $\frac{1 \cdot 2}{(1-3)(2-3)} = 1$.
So $S = 3 \cdot 1 - 3 \cdot 12 + 1 \cdot 10 = -23 \equiv 11 \pmod{17}$. (The polynomial was $g(x) = 2x^2 + 5x + 11$.) Two shares give two equations in three unknowns: for every candidate $S'$ there is exactly one parabola through $(0, S')$ and the two shares, so every secret is equally possible.
\end{solution}

\begin{exercise}
A company wants to protect the master key of its encrypted backups so that (i) no single administrator can read the backups, (ii) the key survives the loss of up to two of its five administrators. Propose a scheme, and say what it cannot protect against.
\end{exercise}
\begin{solution}
Shamir with $n = 5$, $t = 3$: each administrator holds one share (ideally on a hardware token or encrypted with a personal key). One administrator alone, or two together, learn nothing; any three recover the key, so losing two administrators is tolerated. It cannot protect against three administrators colluding, nor against compromise of the machine on which the key is reconstructed when it is used. Shares must also be checked periodically: a share that turns out to be lost only at recovery time reduces the margin.
\end{solution}

\begin{exercise}
An employee asks why the company escrows the keys used to encrypt e-mail but refuses to escrow the keys used to sign it. Explain.
\end{exercise}
\begin{solution}
If an encryption key is lost, the data encrypted with it is lost too, so a recoverable copy is valuable. If a signing key is lost, nothing is lost: a new pair is generated and old signatures remain verifiable with the old public key. An escrowed signing key, instead, would let the escrow holder (or whoever compromises the escrow) sign in the employee's name, so signatures would no longer prove that the employee signed: non-repudiation would be lost.
\end{solution}
